\section{Criterios concretos de usabilidad}
\label{sec:usabilidad}

Esta sección consolida los criterios de usabilidad que orientarán el
diseño y la evaluación del simulador. La consolidación se hace en tres
capas complementarias: (i) los cinco atributos clásicos de
\textit{usability} de Nielsen (aprendizaje, eficiencia, memorabilidad,
manejo de errores, satisfacción), reformulados para el dominio de la
entrevista por voz; (ii) un subconjunto de las heurísticas de Nielsen
de evaluación heurística (HE), aplicadas a las pantallas de la sala
virtual y el plan de mejora; y (iii) un conjunto de \textbf{indicadores
verificables} (umbrales numéricos o conformidades normativas) que se
medirán durante el desarrollo y las pruebas con usuarios. Cada
indicador se enlaza con un requerimiento de la sección
\ref{sec:requerimientos} para asegurar trazabilidad.

\subsection{Atributos de usabilidad y traducción al dominio}

\begin{itemize}
  \item \textbf{Aprendizaje (\textit{learnability}).} Un usuario nuevo
        debe poder iniciar su primera sesión completa en menos de dos
        minutos desde el alta, sin recurrir al tutorial guiado. La
        metáfora del \textit{dojo} y la línea de tiempo circular se
        eligen explícitamente para acelerar la formación del modelo
        mental. Indicador: \rnf{10} \textit{(time-to-first-session)}.

  \item \textbf{Eficiencia (\textit{efficiency}).} Tras tres sesiones
        exitosas, el usuario debe poder iniciar una nueva entrevista en
        no más de dos clics o un comando de voz, gracias a las
        plantillas por industria persistidas en su perfil. Las acciones
        durante la sesión están disponibles tanto por puntero como por
        comando de voz reservado, lo que reduce la fricción al gesto
        físico. Indicador: \rf{06}, \rf{19}.

  \item \textbf{Memorabilidad (\textit{memorability}).} La disposición
        espacial del avatar y el aura es invariante entre los modos
        individual, peer mock y panel; lo único que cambia es la
        información presentada. Los \textit{badges} y rangos por
        competencia refuerzan el aprendizaje longitudinal. Indicador:
        \rf{13}, satisfacción reportada en pruebas con usuarios.

  \item \textbf{Manejo de errores (\textit{error handling}).} La
        interfaz debe prevenir errores costosos antes de cometerlos
        (confirmación explícita antes de publicar al \textit{feed},
        botón ``rehacer última respuesta''), recuperarlos sin pérdida
        de progreso (cola de reintentos ante caídas del LLM,
        reconexión TURN ante caída de WebRTC) y comunicarlos en
        lenguaje natural por voz del entrevistador, no como diálogos
        modales técnicos. Indicador: \rnf{04}, \rnf{14}.

  \item \textbf{Satisfacción.} La experiencia debe celebrar el proceso
        formativo y no únicamente el resultado, evitando los
        \textit{dark patterns} de comparación absoluta. La gamificación
        prioriza mejora relativa del usuario respecto a su línea base,
        no \textit{rankings} de pares ($\rf{14}$). La transparencia
        sobre el funcionamiento del análisis automático y sus
        eventuales sesgos ($\rnf{14}$) sostiene la confianza
        longitudinal en la herramienta.
\end{itemize}

\subsection{Heurísticas de Nielsen aplicadas al simulador}

Se aplican las diez heurísticas de Nielsen como rejilla de evaluación
durante la revisión interna del prototipo y como criterios formales
para la auditoría heurística externa antes de la primera prueba con
usuarios. La tabla \ref{tab:heuristicas} indica, para cada heurística,
cómo se materializa en el simulador y en qué pantalla se verifica.

\begin{table}[ht]
\centering
\renewcommand{\arraystretch}{1.25}
\footnotesize
\begin{tabularx}{\linewidth}{@{}p{4.5cm} Y l@{}}
\toprule
\textbf{Heurística} & \textbf{Materialización en el simulador} &
\textbf{Pantalla / RF} \\
\midrule
H1. Visibilidad del estado del sistema &
El aura procedural y la línea de tiempo circular comunican el estado
de la sesión y del análisis sin abrir paneles separados. &
Sala virtual / \rf{05} \\
\addlinespace
H2. Correspondencia entre sistema y mundo real &
Lenguaje neutral propio del proceso de selección laboral; metáfora
del \textit{dojo} familiar para el usuario. &
Sala virtual, plan de mejora \\
\addlinespace
H3. Control y libertad del usuario &
Comandos ``pausa'', ``repetir'', ``terminar'' siempre disponibles;
botón ``rehacer última respuesta'' antes del cierre. &
Sala virtual / \rf{06} \\
\addlinespace
H4. Consistencia y estándares &
Misma disposición espacial en los modos individual, peer y panel;
componentes de aura y \textit{tooltip} reutilizables. &
Toda la aplicación \\
\addlinespace
H5. Prevención de errores &
Confirmación explícita antes de publicar grabaciones al \textit{feed};
\textit{onboarding} de permisos antes de la primera sesión. &
Privacidad, comunidad / \rf{21} \\
\addlinespace
H6. Reconocer mejor que recordar &
\textit{Tooltips} contextuales sobre cada icono del aura con criterio
y nivel de confianza; sugerencias del asistente lateral basadas en
historial. &
Sala virtual, asistente / \rf{15}, \rf{16} \\
\addlinespace
H7. Flexibilidad y eficiencia &
Atajos de voz y de teclado declarados en cada control; plantillas
por industria; navegación 100\,\% por teclado para todo flujo. &
Configuración / \rf{19} \\
\addlinespace
H8. Diseño estético y minimalista &
La sala virtual oculta cualquier información no esencial durante la
sesión; las métricas viven en el aura, no en \textit{dashboards}. &
Sala virtual / \rf{05} \\
\addlinespace
H9. Recuperación de errores en lenguaje claro &
Mensajes del entrevistador en español neutro, sin códigos técnicos;
explicación verbal de caídas de red o del LLM. &
Sala virtual / \rnf{14} \\
\addlinespace
H10. Ayuda y documentación &
Tour guiado al primer ingreso, asistente lateral por demanda,
ayuda inferida por inactividad o reintentos repetidos. &
Toda la aplicación / \rf{15}, \rf{16} \\
\bottomrule
\end{tabularx}
\caption{Aplicación de las diez heurísticas de Nielsen al simulador.}
\label{tab:heuristicas}
\end{table}

\subsection{Indicadores verificables}

Los criterios anteriores se traducen en indicadores con umbrales o
conformidades concretas, evaluables durante el desarrollo (mediante
prueba automatizada o telemetría) y durante las pruebas con usuarios
(mediante observación cronometrada y cuestionarios SUS). La tabla
\ref{tab:indicadores} resume el conjunto.

\begin{table}[ht]
\centering
\renewcommand{\arraystretch}{1.25}
\footnotesize
\begin{tabularx}{\linewidth}{@{}p{3.6cm} Y p{3.0cm} p{2.4cm}@{}}
\toprule
\textbf{Indicador} & \textbf{Definición operativa} &
\textbf{Umbral / objetivo} & \textbf{Trazabilidad} \\
\midrule
Tiempo a primera sesión & Tiempo entre alta del usuario y final de su
primera entrevista. & $\le 2$ minutos & \rnf{10} \\
\addlinespace
Latencia del aura & Retardo extremo a extremo entre evento percibido y
actualización visual del aura. & $\le 500$\,ms (p95) & \rnf{01},
\rf{05} \\
\addlinespace
Latencia de respuesta del LLM & Retardo entre fin de turno del
candidato y comienzo de la voz del entrevistador. & $\le 2$\,s (p95)
& \rnf{02}, \rf{02} \\
\addlinespace
Latencia de subtítulos & Retardo de aparición del subtítulo respecto
al fonema de inicio de la pregunta del LLM. & $\le 300$\,ms (p95) &
\rf{18} \\
\addlinespace
Tasa de finalización de tarea & Porcentaje de usuarios que completan
una entrevista de extremo a extremo sin asistencia humana. &
$\ge 90\%$ tras 30\,s de inducción & Pruebas moderadas \\
\addlinespace
Tasa de error en STT &
Word Error Rate sobre transcripciones del candidato en español neutro.
& $\le 12\%$ (referencial) & \rf{03} \\
\addlinespace
Conformidad WCAG 2.2 AA &
Auditoría \texttt{axe-core} en CI sin violaciones \textit{serious} ni
\textit{critical}. & 0 violaciones & \rnf{08} \\
\addlinespace
Puntaje Lighthouse &
Categoría \textit{Accessibility} ejecutada en CI sobre las cinco rutas
críticas. & $\ge 95$ & \rnf{08} \\
\addlinespace
Conformidad de teclado &
Cobertura E2E \textit{keyboard-only} de los flujos de usuario. &
$100\%$ de flujos críticos & \rf{19}, \rnf{08} \\
\addlinespace
SUS (System Usability Scale) &
Cuestionario aplicado a una muestra de candidatos primarios al cierre
de la prueba con usuarios. & $\ge 75$ & Prueba final con
usuarios \\
\addlinespace
Costo de IA por sesión &
Costo medio en USD por sesión completa de 30 minutos. &
$\le \$0{,}20$ & \rnf{13} \\
\bottomrule
\end{tabularx}
\caption{Indicadores verificables de usabilidad y conformidad,
con umbrales operativos y trazabilidad a requerimientos.}
\label{tab:indicadores}
\end{table}

\subsection{Estrategia de evaluación}

La verificación de los indicadores anteriores se distribuye en cuatro
hitos a lo largo del ciclo de desarrollo: (i) revisión heurística
interna por dos miembros del equipo siguiendo la rejilla de la tabla
\ref{tab:heuristicas} antes de la primera prueba con usuarios,
(ii) ejecución continua de \texttt{axe-core} y \textit{Lighthouse} en
\textit{integración continua} sobre cada \textit{pull request} para
los indicadores de accesibilidad y rendimiento, (iii) pruebas
moderadas con un usuario representante por perfil primario
(\sk{1}--\sk{7}) sobre una entrevista completa, registrando los
indicadores conductuales (tasa de finalización, tiempo a primera
sesión) y aplicando el cuestionario SUS al cierre, y
(iv) auditoría externa de accesibilidad previa a la entrega final, con
informe firmado de conformidad WCAG 2.2 nivel AA.

Los criterios anteriores no son una capa decorativa: cada uno se
declara como entrada al \textit{backlog} de pruebas automatizadas o al
plan de pruebas con usuarios, de modo que un fallo medible en
cualquiera de ellos bloquea la aceptación del módulo correspondiente.
